\documentclass[10pt,a4paper]{amsart}
\usepackage[margin=19mm]{geometry}
\usepackage{amsmath,amssymb,mathtools}
\usepackage[scr=rsfs,cal=euler]{mathalfa}
\usepackage{xfrac}
\usepackage[unicode,hidelinks,bookmarks=false]{hyperref}
\newcommand{\ie}{i.e.\ }
\newcommand{\cf}{cf.\ }
\newcommand{\resp}{respectively\ }
\newcommand{\Subsection}{\subsection}
\providecommand{\nobreakdash}{\nobreak\hbox{-}}
\setlength{\emergencystretch}{2em}
\multlinegap0pt
\usepackage{xcolor}
\usepackage{tikz}
\usepackage{mathtools}
\usepackage{amssymb}
\usepackage[all]{xy}
\newtheorem{lem}{Lemma}
 \newcommand{\miw}[1]{v_{#1}}
\newcommand{\piv}{\varpi^\lor}
\title{Seidel product formula in equivariant quantum $K$-theory of flag varieties}
\author{Takeshi Ikeda, Takafumi Kouno, Satoshi Naito}
\date{Source: arXiv:2602.24050v1 (CC BY 4.0)}
\begin{document}
\maketitle

\noindent\emph{Source and licence.} Seidel product formula in equivariant quantum $K$-theory of flag varieties. Version arXiv:2602.24050v1 (CC BY 4.0), distributed by arXiv under the Creative Commons Attribution 4.0 International licence, http://creativecommons.org/licenses/by/4.0/, as recorded in the arXiv licence metadata for this exact version and shown on the abstract page https://arxiv.org/abs/2602.24050v1. Full licence notice: \texttt{licenses/arxiv-2602.24050-CC-BY-4.0.txt}.\par\smallskip

\noindent\textit{Editorial scope.} Counted unit: the article's lem lem:key (source0.tex lines 1063--1070). The article's complete original proof of it is delivered with it (source0.tex lines 1073--1083, one \texttt{proof} environment of 11 lines). No mathematics is written, repaired or re-derived: the statement and the proof are the article's own lines, in the article's own notation, and no retained line is substituted or re-flowed.  Retained uncounted as the source-local context the proof needs: lines 684--687; lines 947--951.  Nothing was omitted from any retained span, so the package carries no omission record.  No retained line contains a citation, so the wrapper carries no bibliography and no bibliography entry.  No retained line contains a figure environment, an image-inclusion command or drawing code, so no image is delivered and the package declares no asset dependency.  Editorial formatting only, in addition: an AMS \texttt{amsart} preamble that restates the article's own declarations and macros used by the retained lines, namely the environments \texttt{lem} on separate counters, together with the standard packages those lines need.  The retained environments print in this document as Lemma 1 in order of appearance, whereas the article prints the same items with the numbering of its own full document.  The complete native source of the article as distributed in the arXiv e-print for this exact version is bundled unchanged as \texttt{sources/195442/source0.tex}.
\begin{align}
\miw{i}t_{-\varpi_i^\lor}&=\pi_i^{-1},\label{eq:pi inverse}\\
\pi_i\miw{i}^{-1}\pi_i^{-1}&=\kappa_i\label{eq:min-i and kappa}.
\end{align} 

\begin{lem}\label{lem:key key} For $w\in W$ and $i\in I^s$, we have
\begin{equation}
\gamma_{\miw{i} w}=\gamma_w-w^{-1}(\piv_i).\label{eq:key}
\end{equation}
\end{lem}

\begin{lem}
\label{lem:key}
Let $i$ be any special node, and $w\in W.$ We have
\begin{equation*}
\pi_i^{-1}wt_{\gamma_w}=
\miw{i} wt_{\gamma_{\miw{i} w}}.
\end{equation*}
\end{lem}

\begin{proof}From \eqref{eq:pi inverse}, we have
\begin{align*}
\pi_i^{-1}w
t_{\gamma_w}=
\miw{i} t_{-\piv_i}w
t_{\gamma_w}
=\miw{i}
wt_{-w^{-1}(\piv_i)}t_{\gamma_w}.
\end{align*}
Thus the lemma follows from Lemma \ref{lem:key key}.
\end{proof}

\end{document}
